\documentclass[12pt]{article}
\usepackage{latexblog}

% ============================================================
% Blog Metadata
% ============================================================
\blogtitle{LaTeX(2)：插入图片}
\blogdate{2020-02-08}
\blogtags{TeX, ~LaTeX入门, LaTeX, 其他}
\bloglang{zh}
\blogsummary{}

\begin{document}
\makeblogtitle

在LaTeX中插入图片有些类似于Markdown中的方式。

\section{Figure和图片}\label{figureux548cux56feux7247}

可以使用Figure的方式插入图片：

\begin{Shaded}
\begin{Highlighting}[]
\BuiltInTok{\textbackslash{}usepackage}\NormalTok{\{}\ExtensionTok{graphicx}\NormalTok{\} }\CommentTok{\% 需要导入graphicx包}

\KeywordTok{\textbackslash{}begin}\NormalTok{\{}\ExtensionTok{figure}\NormalTok{\}[h!]}
    \BuiltInTok{\textbackslash{}includegraphics}\NormalTok{[width=}\FunctionTok{\textbackslash{}linewidth}\NormalTok{]\{}\ExtensionTok{src/1.jpg}\NormalTok{\}}
    \FunctionTok{\textbackslash{}caption}\NormalTok{\{冠状病毒的示意图\}}
    \KeywordTok{\textbackslash{}label}\NormalTok{\{}\ExtensionTok{fig:virus}\NormalTok{\}}
\KeywordTok{\textbackslash{}end}\NormalTok{\{}\ExtensionTok{figure}\NormalTok{\}}
\end{Highlighting}
\end{Shaded}

这样生成的效果如图所示：

值得注意的是，在LaTeX中，图片和表格等是浮动元素，Tex会自动帮你找个地方放着。据说这样做的原因是因为它希望你专注在写作内容上面，等你完成了内容再回头来调整格式。因为没有这种经验，所以不确定这样是否是最好的方式，不过我们可能更习惯直接放在我们想放的位置上，所以这里是必须的：

\begin{Shaded}
\begin{Highlighting}[]
\NormalTok{begin\{figure\}[h]}
\end{Highlighting}
\end{Shaded}

这里后面有一个\texttt{h!}，有这样一些可用的参数：

\begin{itemize}
\tightlist
\item
  \(h\)ere: 就放在这里
\item
  \(t\)op： 放在页面顶部
\item
  \(b\)ottom: 页面底部
\item
  \(p\)age: 放在单独的页面
\item
  \(!\)(Override)：强制放在指定的位置
\end{itemize}

虽然我们已经用了\texttt{h!}来强制将图片放在页面的位置，但是很有可能页面余下的空间并不足以放入这张图片，这个时候，这个选项会失效。如果依然确定要将图片放入到下面，可以引入\texttt{float}包，使用\texttt{H}选项：

\begin{Shaded}
\begin{Highlighting}[]
\BuiltInTok{\textbackslash{}usepackage}\NormalTok{\{}\ExtensionTok{float}\NormalTok{\}}

\KeywordTok{\textbackslash{}begin}\NormalTok{\{}\ExtensionTok{figure}\NormalTok{\}[H]}
\end{Highlighting}
\end{Shaded}

就像下面这样：

\section{多张图片排版}\label{ux591aux5f20ux56feux7247ux6392ux7248}

如果要将多张图片排版到一起，可以有多种方法。

\subsection{\texorpdfstring{使用\texttt{subfig}包}{使用subfig包}}\label{ux4f7fux7528subfigux5305}

\href{https://ctan.org/pkg/subfig}{subfig}是一个可以支持嵌套Figure或者表格的包，这个包是\texttt{subfigure}包的替代品，后者已经被弃用了。

\begin{Shaded}
\begin{Highlighting}[]
\BuiltInTok{\textbackslash{}usepackage}\NormalTok{\{}\ExtensionTok{subfig}\NormalTok{\}}

\KeywordTok{\textbackslash{}begin}\NormalTok{\{}\ExtensionTok{figure}\NormalTok{\}}
    \FunctionTok{\textbackslash{}centering}
    \FunctionTok{\textbackslash{}subfloat}\NormalTok{[信息锅]\{}
        \BuiltInTok{\textbackslash{}includegraphics}\NormalTok{[width=0.4}\FunctionTok{\textbackslash{}linewidth}\NormalTok{, height=80pt]\{}\ExtensionTok{src/3.jpg}\NormalTok{\}}
\NormalTok{    \}}
    \FunctionTok{\textbackslash{}subfloat}\NormalTok{[硬气功]\{}
        \BuiltInTok{\textbackslash{}includegraphics}\NormalTok{[width=0.4}\FunctionTok{\textbackslash{}linewidth}\NormalTok{, height=80pt]\{}\ExtensionTok{src/4.jpg}\NormalTok{\}}
\NormalTok{    \}}
    \FunctionTok{\textbackslash{}caption}\NormalTok{\{两种气功流派\}}
    \KeywordTok{\textbackslash{}label}\NormalTok{\{}\ExtensionTok{fig:qg}\NormalTok{\}}
\KeywordTok{\textbackslash{}end}\NormalTok{\{}\ExtensionTok{figure}\NormalTok{\}}
\end{Highlighting}
\end{Shaded}

这样排版出来的效果如下：

\subsection{\texorpdfstring{使用\texttt{subfloat}包}{使用subfloat包}}\label{ux4f7fux7528subfloatux5305}

\href{www.ctan.org/pkg/subfloat}{subfloat}是另一个包，用法如下：

\begin{Shaded}
\begin{Highlighting}[]
\KeywordTok{\textbackslash{}begin}\NormalTok{\{}\ExtensionTok{subfigures}\NormalTok{\}}
\KeywordTok{\textbackslash{}begin}\NormalTok{\{}\ExtensionTok{figure}\NormalTok{\}}
    \FunctionTok{\textbackslash{}centering}
    \FunctionTok{\textbackslash{}fbox}\NormalTok{\{}
        \BuiltInTok{\textbackslash{}includegraphics}\NormalTok{[width=0.4}\FunctionTok{\textbackslash{}linewidth}\NormalTok{, height=80pt]\{}\ExtensionTok{src/3.jpg}\NormalTok{\}}
\NormalTok{    \}}
    \FunctionTok{\textbackslash{}caption}\NormalTok{\{信息锅\}}
\KeywordTok{\textbackslash{}end}\NormalTok{\{}\ExtensionTok{figure}\NormalTok{\}}
\KeywordTok{\textbackslash{}begin}\NormalTok{\{}\ExtensionTok{figure}\NormalTok{\}}
    \FunctionTok{\textbackslash{}centering}
    \FunctionTok{\textbackslash{}fbox}\NormalTok{\{}
        \BuiltInTok{\textbackslash{}includegraphics}\NormalTok{[width=0.4}\FunctionTok{\textbackslash{}linewidth}\NormalTok{, height=80pt]\{}\ExtensionTok{src/4.jpg}\NormalTok{\}}
\NormalTok{    \}}
    \FunctionTok{\textbackslash{}caption}\NormalTok{\{硬气功\}}
\KeywordTok{\textbackslash{}end}\NormalTok{\{}\ExtensionTok{figure}\NormalTok{\}}
\KeywordTok{\textbackslash{}end}\NormalTok{\{}\ExtensionTok{subfigures}\NormalTok{\}}
\end{Highlighting}
\end{Shaded}

这样出来的效果是这样：

两张图片不在一行，没有找到怎么把他们放到一起的方法，也不知道是否支持。

\subsection{\texorpdfstring{使用\texttt{subcaption}包}{使用subcaption包}}\label{ux4f7fux7528subcaptionux5305}

另有使用\href{http://www.ctan.org/pkg/caption}{caption}和\href{http://www.ctan.org/pkg/subcaption}{subcaption}包结合的方式。\texttt{subcaption}包更加新一些，可能更适合选用。

\begin{Shaded}
\begin{Highlighting}[]
\BuiltInTok{\textbackslash{}usepackage}\NormalTok{\{}\ExtensionTok{caption}\NormalTok{\}}
\BuiltInTok{\textbackslash{}usepackage}\NormalTok{\{}\ExtensionTok{subcaption}\NormalTok{\} }

\KeywordTok{\textbackslash{}begin}\NormalTok{\{}\ExtensionTok{figure}\NormalTok{\}}
\FunctionTok{\textbackslash{}centering}
\FunctionTok{\textbackslash{}subcaptionbox}\NormalTok{\{信息锅\}\{}
    \BuiltInTok{\textbackslash{}includegraphics}\NormalTok{[width=0.40}\FunctionTok{\textbackslash{}textwidth}\NormalTok{, height=80pt]\{}\ExtensionTok{src/3.jpg}\NormalTok{\}}
\NormalTok{\}}
\FunctionTok{\textbackslash{}hfill}
\FunctionTok{\textbackslash{}subcaptionbox}\NormalTok{\{硬气功\}\{}
    \BuiltInTok{\textbackslash{}includegraphics}\NormalTok{[width=0.40}\FunctionTok{\textbackslash{}textwidth}\NormalTok{, height=80pt]\{}\ExtensionTok{src/4.jpg}\NormalTok{\}}
\NormalTok{\}}
\FunctionTok{\textbackslash{}caption}\NormalTok{\{两种气功流派\}}
\KeywordTok{\textbackslash{}end}\NormalTok{\{}\ExtensionTok{figure}\NormalTok{\}}
\end{Highlighting}
\end{Shaded}

注意subcaption包貌似和前面两种方法中的某一种可能不兼容。效果如下：

如果去掉\texttt{\textbackslash{}hfill}之后，效果基本上和第一种差不多了。

参考：

\begin{itemize}
\tightlist
\item
  \href{https://tex.stackexchange.com/questions/16207/image-from-includegraphics-showing-up-in-wrong-location}{Image from \textbackslash includegraphics showing up in wrong location?}
\item
  \href{https://tex.stackexchange.com/questions/122314/figures-what-is-the-difference-between-using-subfig-or-subfigure}{Figures: What is the difference between using subfig or subfigure}
\end{itemize}


\end{document}
